\chapter{Standard conjecture of Hodge type}

This chapter is a faithful transcription of \S3 of Ancona's paper. We recall the
two classical standard conjectures of K\"unneth and Lefschetz type (due to
Grothendieck), record the standing assumption that \(X\) satisfies them
(verified in particular by abelian varieties), construct the motivic and
primitive pairings, and state the standard conjecture of Hodge type. We then
prove the first reduction steps: the kernel/perfectness results, the classical
divisor case, the equivalent codimension-two signature formulation for
fourfolds, and the reduction to a finite field. Finally we state the project's
main theorem (\cref{thm:scht}) and record the forward-referencing proof sketch
assembled in \S8 of the paper. Every downstream result of the project ultimately
serves \cref{thm:scht}.

Throughout the chapter, \(k\) is a base field, \(H^*\) is a fixed classical Weil
cohomology with associated realization functor \(R\) (see the Conventions
chapter), \(X\) is a variety over \(k\) of dimension \(d\), and \(L\) is an ample
divisor on \(X\). We write \(\hh(X)\) for the homological motive of \(X\) with
respect to \(H^*\). The whole chapter is written under
\cref{ass:kun-lef}.

\begin{conjecture}[Standard conjecture of K\"unneth type]
  \label{conj:kunneth}
  \lean{MR4199442StandardConjecturesForAbelianFourfolds.KunnethConjecture}
  % SOURCE: [Ancona], §3, Conjecture (conjkun) (read from references/source/fourfolds_2020.tex)
  % SOURCE QUOTE: "There exists a decomposition
  % \[\mathfrak{h}(X)=\bigoplus_{n=0}^{2d} \mathfrak{h}^n(X)\]
  % such that $R(\mathfrak{h}^n(X))=H^n(X)$ and $h^n(X)=h^{2d-n}(X)^{\vee}(-d).$"
  \textit{Source: Ancona, arXiv:1806.03216, §3.}
  There exists a decomposition
  \[\hh(X) = \bigoplus_{n=0}^{2d} \hh^n(X)\]
  in homological motives such that \(R(\hh^n(X)) = H^n(X)\) for every \(n\), and
  such that the duality
  \[\hh^n(X) = \hh^{2d-n}(X)^{\vee}(-d)\]
  holds. This is a classical conjecture and is assumed in what follows.
\end{conjecture}

\begin{conjecture}[Standard conjecture of Lefschetz type]
  \label{conj:lefschetz}
  \lean{MR4199442StandardConjecturesForAbelianFourfolds.LefschetzConjecture}
  % SOURCE: [Ancona], §3, Conjecture (conjlef) (read from references/source/fourfolds_2020.tex)
  % SOURCE QUOTE: "For all $n \leq d$, there is an isomorphism
  % \[\mathfrak{h}^n(X) \isocan \mathfrak{h}^{2d-n}(X)(d-n)\]
  % induced by the cup product with $L^{d-n}$. Moreover, the K\"unneth decomposition can be refined to a  decomposition
  % \[\mathfrak{h}^n(X)= \mathfrak{h}^{n,\prim}(X)\oplus \mathfrak{h}^{n-2}(X)(-1)\]
  %  which realizes to the primitive decomposition of $H^n(X)$."
  \textit{Source: Ancona, arXiv:1806.03216, §3.}
  For all \(n \leq d\), cup product with \(L^{d-n}\) induces a canonical
  isomorphism
  \[\hh^n(X) \cong \hh^{2d-n}(X)(d-n).\]
  Moreover, the K\"unneth decomposition of \cref{conj:kunneth} refines to a
  primitive Lefschetz decomposition
  \[\hh^n(X) = \hh^{n,\prim}(X) \oplus \hh^{n-2}(X)(-1)\]
  which realizes to the primitive decomposition of \(H^n(X)\). This is a
  classical conjecture and is assumed in what follows.
\end{conjecture}

\begin{remark}
  By the work of Lieberman and Kleiman the conjectures of K\"unneth and Lefschetz
  type are known for abelian varieties.
\end{remark}

\begin{definition}[Assumption: K\"unneth and Lefschetz]
  \label{ass:kun-lef}
  \lean{MR4199442StandardConjecturesForAbelianFourfolds.KunnethLefschetzAssumption}
  \uses{conj:kunneth, conj:lefschetz}
  % SOURCE: [Ancona], §3, Assumption (kun+lef) (read from references/source/fourfolds_2020.tex)
  % SOURCE QUOTE: "From now on we will assume that $X$ verifies the   Conjectures \ref{conjkun} and \ref{conjlef}."
  \textit{Source: Ancona, arXiv:1806.03216, §3.}
  We say that \(X\) satisfies the \emph{K\"unneth--Lefschetz assumption} if it
  verifies both the standard conjecture of K\"unneth type (\cref{conj:kunneth})
  and the standard conjecture of Lefschetz type (\cref{conj:lefschetz}). This
  assumption is satisfied in particular by abelian varieties, and is in force
  throughout the chapter.
\end{definition}

\begin{remark}
  Under \cref{ass:kun-lef} the K\"unneth and Lefschetz isomorphisms combine into
  a canonical isomorphism \(\hh^n(X) \cong \hh^{n}(X)^{\vee}(-n)\), which by
  adjunction yields a pairing
  \[q_n : \hh^n(X) \otimes \hh^{n}(X) \longrightarrow \one(-n),\]
  with respect to which the Lefschetz decomposition
  \(\hh^{n}(X) = \hh^{n,\prim}(X) \oplus \hh^{n-2}(X)(-1)\) is orthogonal.
\end{remark}

\begin{definition}[Motivic pairing]
  \label{def:pairing-mot}
  \lean{MR4199442StandardConjecturesForAbelianFourfolds.motivicPairing}
  \uses{ass:kun-lef, def:chow-motives}
  % SOURCE: [Ancona], §3, Remark and Definition (defpairingmot) (read from references/source/fourfolds_2020.tex)
  % SOURCE QUOTE: "Conjectures \ref{conjkun} and \ref{conjlef} induce an isomorphism
  % \[\mathfrak{h}^n(X) \isocan \mathfrak{h}^{n}(X)^{\vee}(-n).\]
  % By adjunction this gives a pairing
  % \[q_n : \mathfrak{h}^n(X) \otimes \mathfrak{h}^{n}(X)  \longrightarrow \one(-n).\]
  % By construction  the Lefschetz decomposition
  % \[\mathfrak{h}^{n }(X)=\mathfrak{h}^{n,\prim}(X) \oplus \mathfrak{h}^{n-2}(X)(-1)\]
  % is orthogonal with respect to this pairing."
  % SOURCE QUOTE: "For $n\leq d$, we define the pairing
  % \[\langle \cdot, \cdot \rangle_{n,\mot} : \mathfrak{h}^n(X) \otimes \mathfrak{h}^{n}(X)  \longrightarrow \one(-n)\]
  % recursively on $n$, by slightly modifying the pairing $q_n$ of the above remark. We impose that  the Lefschetz decomposition
  % \[\mathfrak{h}^{n }(X)=\mathfrak{h}^{n,\prim}(X) \oplus \mathfrak{h}^{n-2}(X)(-1)\]
  % is still orthogonal, we impose the equality $\langle \cdot, \cdot \rangle_{n,\mot} =\langle \cdot, \cdot \rangle_{n-2,\mot}(-2) $ on $\mathfrak{h}^{n-2}(X)(-1)$  and finally on $\mathfrak{h}^{n,\prim}(X)$ we define \[\langle \cdot, \cdot \rangle_{n,\mot}  = (-1)^{n(n+1)/2}q_n. \]"
  \textit{Source: Ancona, arXiv:1806.03216, §3.}
  For \(n \leq d\) the \emph{motivic pairing}
  \[\langle \cdot, \cdot \rangle_{n,\mathrm{mot}} : \hh^n(X) \otimes \hh^{n}(X)
    \longrightarrow \one(-n)\]
  is defined recursively on \(n\) by modifying the adjunction pairing \(q_n\) of
  the previous remark, as follows:
  \begin{itemize}
    \item the Lefschetz decomposition
      \(\hh^{n}(X) = \hh^{n,\prim}(X) \oplus \hh^{n-2}(X)(-1)\) is required to be
      orthogonal;
    \item on the summand \(\hh^{n-2}(X)(-1)\) one imposes the equality
      \(\langle \cdot, \cdot \rangle_{n,\mathrm{mot}} =
      \langle \cdot, \cdot \rangle_{n-2,\mathrm{mot}}(-2)\);
    \item on the primitive part \(\hh^{n,\prim}(X)\) one sets
      \(\langle \cdot, \cdot \rangle_{n,\mathrm{mot}} = (-1)^{n(n+1)/2} q_n\).
  \end{itemize}
\end{definition}

\begin{remark}\label{rem:hodge-type-pola}
  If \(k\) is embedded in \(\C\), the Betti realization of
  \(\langle \cdot, \cdot \rangle_{n,\mathrm{mot}}\) is a polarization of the Hodge
  structure \(H^n(X)\).
\end{remark}

\begin{definition}[Primitive pairing on numerical cycles]
  \label{def:primitive-pairing}
  \lean{MR4199442StandardConjecturesForAbelianFourfolds.primitivePairing}
  \uses{def:pairing-mot, def:num-cycles}
  % SOURCE: [Ancona], §3, Definition (defpairing) (read from references/source/fourfolds_2020.tex)
  % SOURCE QUOTE: "For $n \leq d/2$, we  define the pairing
  % \[\langle \cdot, \cdot \rangle_{n} : \mathcal{Z}^{n}_{\hom}(X)_{\Q}  \times  \mathcal{Z}^{n}_{\hom}(X)_{\Q}  \longrightarrow \Q=\End(\one)\]
  % as follows. For $\alpha, \beta  \in \mathcal{Z}^{n}_{\hom}(X)_{\Q}  = \Hom (\one, \mathfrak{h}^{2n}(X)(n))$ consider
  % \[(\alpha \otimes \beta)  \in   \Hom (\one, \mathfrak{h}^{2n}(X) \otimes \mathfrak{h}^{2n}(X)(2n))\] and define
  % \[\langle \alpha, \beta \rangle_{n} =      \langle \cdot, \cdot \rangle_{2n,\mot} (2n) \circ  (\alpha \otimes \beta).\]"
  % SOURCE QUOTE: "For $n \leq d/2$, define  $\mathcal{Z}^{n,\prim}_{\num}(X)_{\Q} \subset \mathcal{Z}^{n}_{\num}(X)_{\Q}$ as
  % \[\mathcal{Z}^{n,\prim}_{\num}(X)_{\Q}  = \Hom_{\NUM(k)_\Q}(\one, \mathfrak{h}^{2n,\prim}(X)(n)).\]  We will keep the same notation as in (1) for the induced pairing
  % \[\langle \cdot, \cdot \rangle_{n} : \mathcal{Z}^{n,\prim}_{\num}(X)_{\Q}  \times  \mathcal{Z}^{n,\prim}_{\num}(X)_{\Q}  \longrightarrow \Q.\]"
  % SOURCE QUOTE: "The  above definition is equivalent to the one given in the introduction, namely
  % \[\mathcal{Z}^{n,\prim}_{\num}(X)_{\Q} = \{\alpha \in \mathcal{Z}^{n}_{\num}(X)_{\Q} , \hspace{0.2cm} \alpha \cdot L^ {{d-2n+1}}=0 \hspace{0.2cm}  \textrm{in}  \hspace{0.2cm}  \mathcal{Z}^{d-n+1}_{\num}(X)_{\Q}\}.\]
  % Moreover,  for $\alpha$ and $\beta$ in $\mathcal{Z}^{n,\prim}_{\num}(X)_{\Q}$,  we have the equality \[\langle \alpha, \beta \rangle_{n} =  (-1)^n \alpha \cdot \beta \cdot L^{d-2n}.\]"
  \textit{Source: Ancona, arXiv:1806.03216, §3.}
  Fix \(n \leq d/2\). For
  \(\alpha, \beta \in \mathcal{Z}^{n}_{\hom}(X)_{\Q} =
  \Hom(\one, \hh^{2n}(X)(n))\), form
  \(\alpha \otimes \beta \in
  \Hom(\one, \hh^{2n}(X) \otimes \hh^{2n}(X)(2n))\) and define
  \[\langle \alpha, \beta \rangle_{n} =
    \langle \cdot, \cdot \rangle_{2n,\mathrm{mot}}(2n) \circ (\alpha \otimes \beta)
    \;\in\; \Q = \End(\one).\]
  The \emph{space of primitive numerical cycles} is the subspace
  \[\mathcal{Z}^{n,\prim}_{\num}(X)_{\Q} =
    \Hom_{\NUM(k)_\Q}\bigl(\one, \hh^{2n,\prim}(X)(n)\bigr)
    \subset \mathcal{Z}^{n}_{\num}(X)_{\Q},\]
  on which \(\langle \cdot, \cdot \rangle_{n}\) induces a pairing
  \[\langle \cdot, \cdot \rangle_{n} :
    \mathcal{Z}^{n,\prim}_{\num}(X)_{\Q} \times
    \mathcal{Z}^{n,\prim}_{\num}(X)_{\Q} \longrightarrow \Q.\]
  This admits the equivalent elementary description
  \[\mathcal{Z}^{n,\prim}_{\num}(X)_{\Q} =
    \{\alpha \in \mathcal{Z}^{n}_{\num}(X)_{\Q} :
    \alpha \cdot L^{d-2n+1} = 0 \text{ in } \mathcal{Z}^{d-n+1}_{\num}(X)_{\Q}\},\]
  and for \(\alpha, \beta \in \mathcal{Z}^{n,\prim}_{\num}(X)_{\Q}\) one has
  \[\langle \alpha, \beta \rangle_{n} = (-1)^n\, \alpha \cdot \beta \cdot L^{d-2n}.\]
\end{definition}

\begin{conjecture}[Standard conjecture of Hodge type]
  \label{conj:hodge-type}
  \lean{MR4199442StandardConjecturesForAbelianFourfolds.HodgeTypeConjecture}
  \uses{def:primitive-pairing}
  % SOURCE: [Ancona], §3, Conjecture (conjscht) (read from references/source/fourfolds_2020.tex)
  % SOURCE QUOTE: "For all $n \leq d/2$, the pairing  \[\langle \cdot, \cdot \rangle_{n} : \mathcal{Z}^{n,\prim}_{\num}(X)_{\Q}  \times  \mathcal{Z}^{n,\prim}_{\num}(X)_{\Q}  \longrightarrow \Q\]
  % of Definition \ref{defpairing} is positive definite."
  \textit{Source: Ancona, arXiv:1806.03216, §3.}
  For all \(n \leq d/2\), the pairing
  \[\langle \cdot, \cdot \rangle_{n} :
    \mathcal{Z}^{n,\prim}_{\num}(X)_{\Q} \times
    \mathcal{Z}^{n,\prim}_{\num}(X)_{\Q} \longrightarrow \Q\]
  of \cref{def:primitive-pairing} is positive definite.
\end{conjecture}

\begin{proposition}[Kernel of the pairing]
  \label{prop:ker-pairing}
  \lean{MR4199442StandardConjecturesForAbelianFourfolds.kernel_pairing_eq_num_trivial}
  \uses{def:primitive-pairing, def:num-equiv}
  % SOURCE: [Ancona], §3, Proposition (kerpairing) (read from references/source/fourfolds_2020.tex)
  % SOURCE QUOTE: "The kernel of the pairing $\langle \cdot, \cdot \rangle_{n}$  on the vector space $\mathcal{Z}_{\hom}^{n}(X)_{\Q}$ is the space of algebraic classes  which are numerically trivial. The analogous statement holds true for $\mathcal{Z}_{\hom}^{n,\prim}(X)_{\Q}$."
  \textit{Source: Ancona, arXiv:1806.03216, §3.}
  The kernel of the pairing \(\langle \cdot, \cdot \rangle_{n}\) on the vector
  space \(\mathcal{Z}_{\hom}^{n}(X)_{\Q}\) is exactly the space of algebraic
  classes that are numerically trivial. The analogous statement holds for
  \(\mathcal{Z}_{\hom}^{n,\prim}(X)_{\Q}\).
\end{proposition}

% SOURCE QUOTE PROOF: "This is a direct consequence of a general fact. Let $T$ be a homological motive and suppose that it is endowed with an isomorphism
% $f:T \isocan T^\vee$. Call $q: T \otimes T \rightarrow \one $ the pairing induced by adjunction. In analogy to Definition \ref{defpairing}, $q$ induces a pairing $\langle \cdot , \cdot \rangle$ on the space $\Hom(\one, T).$ First, notice that this pairing can also be described in the following way
% \[\langle \alpha,\beta\rangle=\alpha^\vee \circ f \circ \beta.\]
% This equality holds for formal reasons in any rigid category (alternatively, for homological motives, one can check it after realization).
% Using this equality one has that $\beta$ is in the kernel of the pairing if and only if  $(\alpha^\vee \circ f) \circ \beta=0$ for all $\alpha$. As $f$ is an isomorphism, this is equivalent to the fact that $\gamma \circ \beta=0$ for all $\gamma\in\Hom(T,\one).$ This means precisely that $\beta$ is numerically trivial.
% To conclude, notice that, by construction, one can apply this general fact to $T= \mathfrak{h}^{2n}(X)(n)$ or $T= \mathfrak{h}^{2n,\prim}(X)(n)$."
\begin{proof}
  \uses{def:primitive-pairing, def:num-equiv}
  This follows from a general fact. Let \(T\) be a homological motive endowed
  with an isomorphism \(f : T \cong T^\vee\), and let
  \(q : T \otimes T \to \one\) be the pairing it induces by adjunction. As in
  \cref{def:primitive-pairing}, \(q\) induces a pairing
  \(\langle \cdot, \cdot \rangle\) on \(\Hom(\one, T)\), which admits the
  description
  \[\langle \alpha, \beta \rangle = \alpha^\vee \circ f \circ \beta.\]
  This equality holds for formal reasons in any rigid category (and, for
  homological motives, may be checked after realization). Hence \(\beta\) lies in
  the kernel if and only if \((\alpha^\vee \circ f) \circ \beta = 0\) for all
  \(\alpha\); since \(f\) is an isomorphism, this is equivalent to
  \(\gamma \circ \beta = 0\) for all \(\gamma \in \Hom(T, \one)\), i.e. \(\beta\)
  is numerically trivial. Applying this to \(T = \hh^{2n}(X)(n)\) and to
  \(T = \hh^{2n,\prim}(X)(n)\) gives both statements.
\end{proof}

\begin{corollary}[Perfectness and semipositivity]
  \label{cor:semipos}
  \lean{MR4199442StandardConjecturesForAbelianFourfolds.pairing_perfect_and_semipos}
  \uses{prop:ker-pairing}
  % SOURCE: [Ancona], §3, Corollary (semipos) (read from references/source/fourfolds_2020.tex)
  % SOURCE QUOTE: "The pairing $\langle \cdot, \cdot \rangle_{n}$  is perfect on $\mathcal{Z}_{\num}^{n}(X)_{\Q}$. Moreover the pairing  is positive definite on  $\mathcal{Z}_{\num}^{n}(X)_{\Q}$ if and only if it is positive semidefinite on $\mathcal{Z}_{\hom}^{n}(X)_{\Q}$. The analogous statements hold true for $\mathcal{Z}_{\num}^{\prim,n}(X)_{\Q}$."
  \textit{Source: Ancona, arXiv:1806.03216, §3.}
  The pairing \(\langle \cdot, \cdot \rangle_{n}\) is perfect on
  \(\mathcal{Z}_{\num}^{n}(X)_{\Q}\). Moreover it is positive definite on
  \(\mathcal{Z}_{\num}^{n}(X)_{\Q}\) if and only if it is positive semidefinite
  on \(\mathcal{Z}_{\hom}^{n}(X)_{\Q}\). The analogous statements hold for
  \(\mathcal{Z}_{\num}^{n,\prim}(X)_{\Q}\).
\end{corollary}

\begin{proof}
  \uses{prop:ker-pairing}
  By \cref{prop:ker-pairing} the kernel of \(\langle \cdot, \cdot \rangle_{n}\)
  on \(\mathcal{Z}_{\hom}^{n}(X)_{\Q}\) is precisely the space of numerically
  trivial classes. The numerical space \(\mathcal{Z}_{\num}^{n}(X)_{\Q}\) is the
  quotient of \(\mathcal{Z}_{\hom}^{n}(X)_{\Q}\) by that kernel, so the induced
  pairing on \(\mathcal{Z}_{\num}^{n}(X)_{\Q}\) has trivial kernel, i.e. is
  perfect. Since a symmetric form that is positive semidefinite descends to a
  positive definite form precisely on the quotient by its kernel, positive
  semidefiniteness on \(\mathcal{Z}_{\hom}^{n}(X)_{\Q}\) is equivalent to
  positive definiteness on \(\mathcal{Z}_{\num}^{n}(X)_{\Q}\). The same argument
  applies verbatim with \(\hh^{2n}(X)(n)\) replaced by \(\hh^{2n,\prim}(X)(n)\),
  giving the primitive statements.
\end{proof}

\begin{theorem}[Divisor case, Segre--Bronowski]
  \label{thm:segre}
  \lean{MR4199442StandardConjecturesForAbelianFourfolds.segre_divisor_positive_definite}
  \uses{def:primitive-pairing}
  % SOURCE: [Ancona], §3, Theorem (segre) (read from references/source/fourfolds_2020.tex)
  % SOURCE QUOTE: "Together with the case of characteristic zero (see the above remark), the following is the only known result on the standard conjecture of Hodge type \cite[5.3.2.3]{Andmot}."
  % SOURCE QUOTE: "The pairing\label{segre} $\langle \cdot, \cdot \rangle_{1}$   on $\mathcal{Z}_{\num}^{1}(X)_{\Q}$ is positive definite. Equivalently, the pairing
  % \[\mathcal{Z}^1_{\num}(X)_{\Q}\times \mathcal{Z}_{\num}^1(X)_{\Q} \longrightarrow \Q \hspace{1cm} D,D' \mapsto D\cdot D' \cdot L^{d-2}\]
  % is of signature  $ (s_+;s_-) = (1;  \rho_1-1)$ with
  %  $\rho_1=\dim \mathcal{Z}_{\num}^1(X)_{\Q}$."
  \textit{Source: Ancona, arXiv:1806.03216, §3.}
  The pairing \(\langle \cdot, \cdot \rangle_{1}\) on
  \(\mathcal{Z}_{\num}^{1}(X)_{\Q}\) is positive definite. Equivalently, the
  pairing
  \[\mathcal{Z}^1_{\num}(X)_{\Q} \times \mathcal{Z}_{\num}^1(X)_{\Q}
    \longrightarrow \Q, \qquad (D, D') \mapsto D \cdot D' \cdot L^{d-2}\]
  has signature \((s_+; s_-) = (1; \rho_1 - 1)\), where
  \(\rho_1 = \dim \mathcal{Z}_{\num}^1(X)_{\Q}\).
\end{theorem}

\begin{proof}
  This is the classical Hodge index theorem for divisors (Segre--Bronowski). It
  is the only previously known case of the standard conjecture of Hodge type in
  positive characteristic, and we cite it as a known input (André,
  \emph{Une introduction aux motifs}, 5.3.2.3). No proof is reproduced here.
\end{proof}

\begin{proposition}[Equivalent signature formulation for fourfolds]
  \label{prop:any-polarization}
  \lean{MR4199442StandardConjecturesForAbelianFourfolds.hodge_type_iff_signature}
  \uses{thm:segre, conj:hodge-type, def:primitive-pairing}
  % SOURCE: [Ancona], §3, Proposition (anypolarization) (read from references/source/fourfolds_2020.tex)
  % SOURCE QUOTE: "The standard conjecture of Hodge type for a fourfold $X$  is equivalent to the statement that the intersection product
  % \[\mathcal{Z}^2_{\num}(X)_{\Q}\times \mathcal{Z}_{\num}^2(X)_{\Q} \longrightarrow \Q\]
  % is of signature  $ (s_+;s_-) = ( \rho_2 - \rho_1+1;  \rho_1-1)$ with
  %  $\rho_n=\dim \mathcal{Z}_{\num}^n(X)_{\Q}$.
  % In particular the conjecture does not depend on the ample divisor $L$ chosen."
  \textit{Source: Ancona, arXiv:1806.03216, §3.}
  For a fourfold \(X\), the standard conjecture of Hodge type
  (\cref{conj:hodge-type}) is equivalent to the assertion that the intersection
  product
  \[\mathcal{Z}^2_{\num}(X)_{\Q} \times \mathcal{Z}_{\num}^2(X)_{\Q}
    \longrightarrow \Q\]
  has signature \((s_+; s_-) = (\rho_2 - \rho_1 + 1; \rho_1 - 1)\), where
  \(\rho_n = \dim \mathcal{Z}_{\num}^n(X)_{\Q}\). In particular the conjecture is
  independent of the chosen ample divisor \(L\).
\end{proposition}

% SOURCE QUOTE PROOF: "By Theorem \ref{segre} the standard conjecture of Hodge type holds true for divisors.  Hence we have to study codimension $2$ cycles.
% Consider the decomposition \[\mathcal{Z}^{2}_{\num}(X)_{\Q} = \mathcal{Z}^{2,\prim}_{\num}(X)_{\Q} \oplus L \cdot \mathcal{Z}^{1}_{\num}(X)_{\Q}\] induced by the Lefschetz decomposition \[\mathfrak{h}^{4}(X)=\mathfrak{h}^{4,\prim}(X) \oplus \mathfrak{h}^{2}(X)(-1).\]  It is orthogonal with respect to the intersection product as the Lefschetz decomposition is orthogonal with respect to the cup product (already at homological level).
% We claim that the intersection product on  $L \cdot \mathcal{Z}^{1}_{\num}(X)_{\Q}$ is of signature $(1;  \rho_1-1)$. Assuming the claim notice that  the intersection product is positive definite on $\mathcal{Z}^{2,\prim}_{\num}(X)_{\Q}$ if and only if the intersection product of the total space $\mathcal{Z}^2_{\num}(X)_{\Q}$ has the predicted signature.
% For the claim, consider $\alpha=L \cdot D$ an element of $L \cdot \mathcal{Z}^{1}_{\num}(X)_{\Q}$. Then we have
% $\alpha\cdot \alpha = D \cdot D \cdot L^2$
% hence the claim follows from Theorem \ref{segre}."
\begin{proof}
  \uses{thm:segre, def:primitive-pairing}
  By \cref{thm:segre} the standard conjecture of Hodge type holds for divisors,
  so it remains to study codimension-two cycles. The Lefschetz decomposition
  \(\hh^{4}(X) = \hh^{4,\prim}(X) \oplus \hh^{2}(X)(-1)\) induces a decomposition
  \[\mathcal{Z}^{2}_{\num}(X)_{\Q} =
    \mathcal{Z}^{2,\prim}_{\num}(X)_{\Q} \oplus L \cdot \mathcal{Z}^{1}_{\num}(X)_{\Q},\]
  orthogonal for the intersection product because the Lefschetz decomposition is
  orthogonal for the cup product already at the homological level.

  We claim the intersection product on \(L \cdot \mathcal{Z}^{1}_{\num}(X)_{\Q}\)
  has signature \((1; \rho_1 - 1)\). Granting the claim, the intersection product
  is positive definite on \(\mathcal{Z}^{2,\prim}_{\num}(X)_{\Q}\) if and only if
  the intersection product on the total space \(\mathcal{Z}^2_{\num}(X)_{\Q}\) has
  the predicted signature \((\rho_2 - \rho_1 + 1; \rho_1 - 1)\); and positive
  definiteness on the primitive part is exactly the standard conjecture of Hodge
  type for \(X\) in codimension two (via \cref{def:primitive-pairing}, noting
  \(\langle \alpha, \beta \rangle_2 = \alpha \cdot \beta\) on the primitive part
  of a fourfold, \(d - 2n = 0\)). For the claim, an element
  \(\alpha = L \cdot D \in L \cdot \mathcal{Z}^{1}_{\num}(X)_{\Q}\) satisfies
  \(\alpha \cdot \alpha = D \cdot D \cdot L^2\), so the claim follows from
  \cref{thm:segre}. Since the formulation refers only to the intersection product
  on \(\mathcal{Z}^2_{\num}(X)_{\Q}\), it does not depend on \(L\).
\end{proof}

\begin{proposition}[Reduction to a finite field]
  \label{prop:reduce-finite}
  \lean{MR4199442StandardConjecturesForAbelianFourfolds.reduce_to_finite_field}
  \uses{cor:semipos, ass:kun-lef}
  % SOURCE: [Ancona], §3, Proposition (reducefinite) (read from references/source/fourfolds_2020.tex)
  % SOURCE QUOTE: "Let $X_0$ be a specialization of $X$. Suppose that $X_0$ as well satisfies the Assumption \ref{kun+lef} (with respect to the specialization of $L$). Then the standard conjecture of Hodge type for $X_0$ implies the standard conjecture of Hodge type for $X$."
  \textit{Source: Ancona, arXiv:1806.03216, §3.}
  Let \(X_0\) be a specialization of \(X\), and suppose \(X_0\) also satisfies
  \cref{ass:kun-lef} (with respect to the specialization of \(L\)). Then the
  standard conjecture of Hodge type for \(X_0\) implies the standard conjecture of
  Hodge type for \(X\). In particular the conjecture reduces to varieties defined
  over a finite field.
\end{proposition}

% SOURCE QUOTE PROOF: "Consider the canonical inclusion  $\mathcal{Z}_{\hom}^{n,\prim}(X)_{\Q}\subset \mathcal{Z}_{\hom}^{n,\prim}(X_0)_{\Q}$ induced by smooth proper base change in $\ell$-adic cohomology. Then, if the pairing $\langle \cdot, \cdot \rangle_{n}$  is positive semidefinite  on $\mathcal{Z}_{\hom}^{n,\prim}(X_0)_{\Q}$ it must be so for  $\mathcal{Z}_{\hom}^{n,\prim}(X)_{\Q} $. On the other hand, by Corollary \ref{semipos}, this semipositivity property is a reformulation of the standard conjecture of Hodge type."
\begin{proof}
  \uses{cor:semipos}
  Smooth proper base change in \(\ell\)-adic cohomology gives a canonical
  inclusion
  \(\mathcal{Z}_{\hom}^{n,\prim}(X)_{\Q} \subset
  \mathcal{Z}_{\hom}^{n,\prim}(X_0)_{\Q}\). Hence if
  \(\langle \cdot, \cdot \rangle_{n}\) is positive semidefinite on
  \(\mathcal{Z}_{\hom}^{n,\prim}(X_0)_{\Q}\) it is positive semidefinite on the
  subspace \(\mathcal{Z}_{\hom}^{n,\prim}(X)_{\Q}\). By \cref{cor:semipos} this
  semipositivity is equivalent to the standard conjecture of Hodge type, so the
  conjecture passes from \(X_0\) to \(X\).
\end{proof}

The following is the main result of the project; it is proved at the end of \S8
of the paper, and the proof sketch is recorded below.

\begin{theorem}[Standard conjecture of Hodge type for abelian fourfolds]
  \label{thm:scht}
  \lean{MR4199442StandardConjecturesForAbelianFourfolds.standardConjectureHodgeType_abelianFourfold}
  \uses{conj:hodge-type, prop:reduce-finite, prop:exotic-dim2, thm:honda-tate, prop:any-polarization, cor:cm-stable, lem:lefschetz-pairings, prop:scht-lefschetz, prop:factor-algebraic, cor:lifting-motives, thm:mainthm}
  % SOURCE: [Ancona], §3, Theorem (thmscht) (read from references/source/fourfolds_2020.tex)
  % SOURCE QUOTE: "The standard conjecture of Hodge type holds for abelian fourfolds in positive characteristic."
  \textit{Source: Ancona, arXiv:1806.03216, §3.}
  The standard conjecture of Hodge type (\cref{conj:hodge-type}) holds for abelian
  fourfolds in positive characteristic.
\end{theorem}

% SOURCE QUOTE PROOF: "Let $A$ be an abelian fourfold and $L$ be a hyperplane section. By Proposition \ref{reducefinite} we can suppose that $A$ and $L$ are defined over a finite field $k=\mathbb{F}_q$. Notice that, in order to prove Theorem \ref{thmscht}, we can (and will) replace $k$ by a finite extension.
% Consider now the decomposition from Proposition \ref{primadecomposizione}(\ref{decomposizioneQ}) and, among the factors of this decomposition, consider
% those that are exotic (Definition \ref{defexotic}). Then there exist a finite extension $k'$ of $k$ and a  CM-structure for $A\times_k k'$ such that any exotic motive  of $\mathfrak{h}^{4}(A\times_k k')$ has dimension two (Proposition \ref{exoticdim2}).  By Theorem \ref{HondaTate}, the CM-structure can be lifted    to characteristic zero.  Moreover, by Proposition \ref{anypolarization}, it is enough to work with a single $L$ for a given abelian fourfold. By Corollary \ref{CMstab} we can choose such an $L$ so that it lifts to characteristic zero and the CM-structure is $*$-stable (Definition \ref{definitionCM}).
% Now, by Lemma \ref{accoppiamenti1} we can work with the pairing $\langle \cdot, \cdot \rangle_{1,\mot}^{\otimes 4}$ instead of $\langle \cdot, \cdot \rangle_{4,\mot}$. For this pairing, the decomposition from Proposition \ref{primadecomposizione}(\ref{decomposizioneQ}) is orthogonal, hence we can work with a single motive of the decomposition. By Propositions \ref{oklefschetz} and \ref{talvez}(\ref{generatolefschetz}) we are reduced to motives that are exotic.
% Finally, those exotic motives $\mathcal{M}_I$ are settled by Theorem \ref{mainthm} by setting $M=\mathcal{M}_I(2)$. Notice that all the hypothesis of Theorem \ref{mainthm} are satisfied. Indeed, the motive lives in mixed characteristic (Corollary \ref{lifting motives}), together with its quadratic form (because $L$ lifts to characteristic zero). The space $V_B$ is clearly of dimension two; so it is $V_Z$ by   Proposition  \ref{talvez}(\ref{generatoalgebrico}).  Last, the quadratic form $q_B=R_B(\langle \cdot, \cdot \rangle_{1,\mot}^{\otimes 4})$ is a polarization as $R_B(\langle \cdot, \cdot \rangle_{1,\mot})$ is so."
\begin{proof}
  \uses{prop:reduce-finite, prop:exotic-dim2, thm:honda-tate, prop:any-polarization, cor:cm-stable, lem:lefschetz-pairings, prop:scht-lefschetz, prop:factor-algebraic, cor:lifting-motives, thm:mainthm}
  Let \(A\) be an abelian fourfold and \(L\) a hyperplane section. By
  \cref{prop:reduce-finite} we may assume \(A\) and \(L\) are defined over a
  finite field \(k = \F_q\); to prove the theorem we are free to replace \(k\) by
  a finite extension.

  Among the factors of the canonical decomposition of \(\hh^4(A)\), consider
  those that are \emph{exotic}. There exist a finite extension \(k'\) of \(k\)
  and a CM-structure for \(A \times_k k'\) such that every exotic motive of
  \(\hh^{4}(A \times_k k')\) has dimension two (\cref{prop:exotic-dim2}). By
  \cref{thm:honda-tate} this CM-structure lifts to characteristic zero. By
  \cref{prop:any-polarization} it is enough to work with a single \(L\) for a
  given abelian fourfold, and by \cref{cor:cm-stable} we may choose \(L\) so that
  it lifts to characteristic zero and the CM-structure is \(*\)-stable.

  By \cref{lem:lefschetz-pairings} we may replace the pairing
  \(\langle \cdot, \cdot \rangle_{4,\mathrm{mot}}\) by
  \(\langle \cdot, \cdot \rangle_{1,\mathrm{mot}}^{\otimes 4}\); for the latter
  the CM-decomposition is orthogonal, so we may work with a single factor of the
  decomposition. By \cref{prop:scht-lefschetz} and \cref{prop:factor-algebraic}
  the Lefschetz factors are disposed of, reducing us to the exotic factors.

  Finally, each exotic factor \(\mathcal{M}_I\) is settled by
  \cref{thm:mainthm} applied with \(M = \mathcal{M}_I(2)\). All hypotheses of
  \cref{thm:mainthm} hold: the motive lives in mixed characteristic
  (\cref{cor:lifting-motives}), together with its quadratic form (because \(L\)
  lifts to characteristic zero); the Betti space is two-dimensional, and so is the
  space of numerical cycles by \cref{prop:factor-algebraic}; and the quadratic
  form \(q_B = R_B(\langle \cdot, \cdot \rangle_{1,\mathrm{mot}}^{\otimes 4})\) is
  a polarization since \(R_B(\langle \cdot, \cdot \rangle_{1,\mathrm{mot}})\) is.
\end{proof}
